% Theorem System
% The following boxes are provided:
%   Definition:     \defn 
%   Theorem:        \thm 
%   Lemma:          \lem
%   Corollary:      \cor
%   Proposition:    \prop   
%   Claim:          \clm
%   Fact:           \fact
%   Proof:          \pf
%   Example:        \ex
%   Remark:         \rmk (sentence), \rmkb (block)
% Suffix
%   r:              Allow Theorem/Definition to be referenced, e.g. thmr
%   p:              Add a short proof block for Lemma, Corollary, Proposition or Claim, e.g. lemp
%                   For theorems, use \pf for proof blocks

\newenvironment{foldblock}{}{}



% --------------------------------------------------------------
%     Definition
% --------------------------------------------------------------
\newtheoremstyle{def}
    {\topsep}
    {\topsep}
    {\normalfont}       % Body font
    {}         
    {\color{ForestGreen}\bfseries}
    {.}
    {0.5em}
    {}
\theoremstyle{def}
\newtheorem{mydefinition}{Definition}[section]

\NewDocumentCommand{\defnn}{+m}{
  \begin{mydefinition}
    #1
  \end{mydefinition}
}

\NewDocumentCommand{\defn}{m+m}{
  \begin{mydefinition}[#1]
    #2
  \end{mydefinition}
}

\NewDocumentCommand{\defnr}{mm+m}{
  \begin{mydefinition}[#1]\label{defn:#2}
    #3
  \end{mydefinition}
}




% --------------------------------------------------------------
%     Theorem
% --------------------------------------------------------------
\newtheoremstyle{thm}
  {\topsep}           % Space above
  {\topsep}           % Space below
  {\normalfont}       % Body font
  {}                  % Indent
  {\color{red}\bfseries} % Head font
  {.}                 % Punctuation after heading
  {0.5em}             % Space after heading
  {} 
\theoremstyle{thm}
\newtheorem{mytheorem}{Theorem}[section]

\NewDocumentCommand{\thm}{m+m}{
    \begin{mytheorem}[#1]{}{}
        \leavevmode
        \begin{adjustwidth}{1.5em}{0pt}
            \setlength{\parskip}{0.3\baselineskip} % 统一段落间距（柔和）
            #2
        \end{adjustwidth}
        \vspace{0.2\baselineskip}
    \end{mytheorem}
}

\NewDocumentCommand{\thmnn}{+m}{
    \begin{mytheorem}{}
        \leavevmode
        \begin{adjustwidth}{1.5em}{0pt}
            \setlength{\parskip}{0.3\baselineskip} % 统一段落间距（柔和）
            #1
        \end{adjustwidth}
        \vspace{0.2\baselineskip}
    \end{mytheorem}
}

\NewDocumentCommand{\thmr}{mm+m}{
    \begin{mytheorem}[#1]{}
        \leavevmode
        \phantomsection        % ← 强制插入 anchor
        \label{thm:#2}         
        \begin{adjustwidth}{1.5em}{0pt}
            \setlength{\parskip}{0.3\baselineskip} % 统一段落间距（柔和）
            #3
        \end{adjustwidth}
        \vspace{0.2\baselineskip}
    \end{mytheorem}
}

\NewDocumentCommand{\thmnnr}{m+m}{
    \begin{mytheorem}{}       
        \phantomsection        % ← 强制插入 anchor
        \label{thm:#1}         
        \begin{adjustwidth}{1.5em}{0pt}
            \setlength{\parskip}{0.3\baselineskip} % 统一段落间距（柔和）
            #2
        \end{adjustwidth}
        \vspace{0.2\baselineskip}
    \end{mytheorem}
}


% --------------------------------------------------------------
%     Lemma
% --------------------------------------------------------------
\newtheoremstyle{lemma}
    {\topsep}
    {\topsep}
    {\normalfont}       % Body font
    {}                    % Indent
    {\color{NavyBlue}\bfseries}
    {.}
    {0.5em}
    {}
\theoremstyle{lemma}
\newtheorem{mylemma}{Lemma}[section]

\NewDocumentCommand{\lem}{m+m}{
    \begin{mylemma}[#1]{}
        \leavevmode
        \begin{adjustwidth}{1.5em}{0pt}
            \setlength{\parskip}{0.3\baselineskip} % 统一段落间距（柔和）
            #2
        \end{adjustwidth}
        \vspace{0.2\baselineskip}
    \end{mylemma}
}

\NewDocumentCommand{\lemnn}{+m}{
    \begin{mylemma}{}
        \leavevmode
        \begin{adjustwidth}{1.5em}{0pt}
            \setlength{\parskip}{0.3\baselineskip} % 统一段落间距（柔和）
            #1
        \end{adjustwidth}
        \vspace{0.2\baselineskip}
    \end{mylemma}
}

\NewDocumentCommand{\lemr}{mm+m}{
    \begin{mylemma}[#1]{\label{lemma:#2}}
        #3
    \end{mylemma}
}


\newenvironment{lempf}{%
  \vspace{3pt}
  \noindent\textbf{\textit{Proof for Lemma.}}\par
  \vspace{5pt}
  \noindent\textcolor{NavyBlue}{\rule[0.5ex]{1ex}{1ex}}.\hspace{0.5em}\ignorespaces
}{
  \hfill\qedsymbol\par\vspace{3pt}
}


\NewDocumentCommand{\lemp}{m+m+m}{
    \begin{mylemma}[#1]{}
        #2
    \end{mylemma}

    \begin{lempf}
        #3
    \end{lempf}
}


% --------------------------------------------------------------
%     Corollary
% --------------------------------------------------------------
\newtheoremstyle{cor}
    {\topsep}
    {\topsep}
    {\normalfont}       % Body font
    {}                   % Indent
    {\color{Violet}\bfseries}
    {.}
    {.5em}
    {}
\theoremstyle{cor}
\newtheorem{mycorollary}{Corollary}[section]

\NewDocumentCommand{\cor}{+m}{
    \begin{mycorollary}{}{}
        #1
    \end{mycorollary}
}

\newenvironment{corpf}{%
  \vspace{3pt}
  \noindent\textbf{\textit{Proof for Corollary.}}\par
  \vspace{5pt}
  \noindent\textcolor{Violet}{\rule[0.5ex]{1ex}{1ex}}.\hspace{0.5em}\ignorespaces
}{
  \hfill\qedsymbol\par\vspace{3pt}
}

\NewDocumentCommand{\ncor}{m+m}{
    \begin{mycorollary}[#1]{}
        \leavevmode
        \begin{adjustwidth}{1.5em}{0pt}
            \setlength{\parskip}{0.3\baselineskip} % 统一段落间距（柔和）
            #2
        \end{adjustwidth}
        \vspace{0.2\baselineskip}
    \end{mycorollary}
}

\NewDocumentCommand{\corp}{m+m+m}{
    \begin{mycorollary}[#1]{}
        \leavevmode
        \begin{adjustwidth}{1.5em}{0pt}
            \setlength{\parskip}{0.3\baselineskip} % 统一段落间距（柔和）
            #2
        \end{adjustwidth}
        \vspace{0.2\baselineskip}
    \end{mycorollary}

    \begin{corpf}
        #3
    \end{corpf}
}

\NewDocumentCommand{\cornnp}{+m+m}{
    \begin{mycorollary}{}
        \leavevmode
        \begin{adjustwidth}{1.5em}{0pt}
            \setlength{\parskip}{0.3\baselineskip} % 统一段落间距（柔和）
            #1
        \end{adjustwidth}
        \vspace{0.2\baselineskip}
    \end{mycorollary}

    \begin{corpf}
        #2
    \end{corpf}
}


% --------------------------------------------------------------
%     Proposition
% --------------------------------------------------------------
\newtheoremstyle{prop}
    {\topsep}
    {\topsep}
    {\normalfont}       % Body font
    {}                     % Indent
    {\color{MidnightBlue}\bfseries}
    {.}
    {.5em}
    {}
\theoremstyle{prop}
\newtheorem{myproposition}{Proposition}[section]


\NewDocumentCommand{\prop}{+m}{
    \begin{myproposition}{}{}
        #1
    \end{myproposition}
}

\newenvironment{proppf}{%
  \vspace{3pt}
  \noindent\textbf{\textit{Proof for Proposition.}}\par
  \vspace{5pt}
  \noindent\textcolor{MidnightBlue}{\rule[0.5ex]{1ex}{1ex}}.\hspace{0.5em}\ignorespaces
}{
  \hfill\qedsymbol\par\vspace{3pt}
}


\NewDocumentCommand{\propp}{+m+m}{
    \begin{myproposition}{}{}
        #1
    \end{myproposition}

    \begin{proppf}
        #2
    \end{proppf}
}

% --------------------------------------------------------------
%     Claim
% --------------------------------------------------------------
\newtheoremstyle{claim}
    {\topsep}
    {\topsep}
    {\normalfont}       % Body font
    {}                     % Indent
    {\color{Blue}\bfseries}
    {.}
    {.5em}
    {}
\theoremstyle{claim}
\newtheorem{myclaim}{Claim}[section]

\NewDocumentCommand{\clm}{m+m}{
    \begin{myclaim}{#1}{}
        #2
    \end{myclaim}
}

\newenvironment{clmpf}{%
  \vspace{3pt}
  \noindent\textbf{\textit{Proof for Claim.}}\par
  \vspace{5pt}
  \noindent\textcolor{Blue}{\rule[0.5ex]{1ex}{1ex}}.\hspace{0.5em}\ignorespaces
}{
  \hfill\qedsymbol\par\vspace{3pt}
}

\NewDocumentCommand{\clmp}{m+m+m}{
    \begin{myclaim}{#1}{}
        #2
    \end{myclaim}

    \begin{clmpf}
        #3
    \end{clmpf}
}

% --------------------------------------------------------------
%     Fact
% --------------------------------------------------------------
\newtheoremstyle{fact}
    {\topsep}
    {\topsep}
    {\normalfont}       % Body font
    {}                    % Indent
    {\color{CadetBlue}\bfseries}
    {.}
    {.5em}
    {}
\theoremstyle{fact}
\newtheorem{myfact}{Fact}[section]

\NewDocumentCommand{\fact}{+m}{
    \begin{myfact}{}{}
        #1
    \end{myfact}
}

\newenvironment{factpf}{%
  \vspace{3pt}
  \noindent\textbf{\textit{Proof for fact.}}\par
  \vspace{5pt}
  \noindent\textcolor{CadetBlue}{\rule[0.5ex]{1ex}{1ex}}.\hspace{0.5em}\ignorespaces
}{
  \hfill\qedsymbol\par\vspace{3pt}
}

\NewDocumentCommand{\factp}{+m}{
    \begin{factpf}
        #1
    \end{factpf}
}


% --------------------------------------------------------------
%     proof
% --------------------------------------------------------------
\NewDocumentCommand{\pf}{+m}{
    \begin{proof}
        [\noindent\textbf{Proof.}]
        #1
    \end{proof}
}

\NewDocumentCommand{\pff}{m+m}{
    \begin{proof}
        [\noindent\textbf{Proof for #1.}]
        #2
    \end{proof}
}



% --------------------------------------------------------------
%     Example
% --------------------------------------------------------------
\newenvironment{example}[2][Example]{\begin{trivlist}
\item[\hskip \labelsep {\textcolor{blue}{\bfseries #1}}\hskip \labelsep {\textcolor{blue}{\bfseries #2}}]}{\end{trivlist}}

% \newenvironment{example}{%
%     \par\vspace{5pt}%
%     \noindent{\textbf{Example.}}\ %
% }{%
%     \par\vspace{5pt}%
% }


\NewDocumentCommand{\ex}{m+m}{
    \begin{example}{#1.}
        #2
    \end{example}
}




% --------------------------------------------------------------
%     Remark
% --------------------------------------------------------------
\NewDocumentCommand{\rmk}{+m}{
    {\it \color{blue!50!white}#1}
}
\newenvironment{remark}{%
    \par\vspace{5pt}%
    \noindent\textcolor{Gray}{\textbf{Remark.}}\ %
}{%
    \par\vspace{5pt}%
}

\NewDocumentCommand{\rmkb}{+m}{
    \begin{remark}
        #1
    \end{remark}
}
